Reflect on and discuss the software engineering practices and processes
adopted this semester. Describe ideas for future work that could improve
the project. Include only practices and techniques actually adopted in
your project. Delete any subsection that does not apply.

% --------------------------------
\section{Scrum}

Reflect on the extent to which Scrum practices improved your project
work. What outcomes and behaviours — e.g.\ better code quality, better
accountability — changed? What challenges were experienced?

% --------------------------------
\section{Requirements Engineering}

The methods the team used for requirements engineering had different effects on the quality of our project. Some were very important during the development process, while others were mostly helpful in the early design stage.

At the start, functional and non-functional requirements were very important because they defined what the system would do, its goals, and helped create the list of tasks to complete. As development moved forward, the process became easier to adjust with new tasks coming up mainly from testing, feedback, and making changes to improve the code.

User stories gave a viewpoint focused on the users of the system. At first, these were part of the course requirements, but they turned out to be helpful for explaining how users work, what they expect, and how they interact with the system. This helped in creating features that are easy to use.

Use case diagrams show a general view of how actors interact with the system. They showed how various users engage with the system's features and explained the connections with outside systems.

High-level requirements served as a middle step that connects broad ideas about the system to more specific functional and non-functional requirements. They were simpler to grasp and offered a clear way to move towards more detailed descriptions.

Activity diagrams were helpful for showing how users interact with the system in a clear, step-by-step format. They gave a straightforward summary of how the system works and helped with the initial design choices.

Sequence diagrams give a clearer picture of how systems interact, highlighting the complete communication between different parts of the system and outside systems. They helped make it clear how information and responsibilities move around within the system.

The UML class diagram had the biggest influence on the project. It acted as the main plan for the building, outlining how the system is organized, the different classes, their duties, and how they are connected. It helped break down the system and share tasks among the team, which made working together better. Even though the class diagram started off as a guide for implementation, it slowly started to stray from the code because of feedback, bug fixes, and changes to improve the code. So, the UML models were changed to show the new way things were being implemented instead of just telling how to do it.

Other methods, like package diagrams, CRC cards, and verb-noun analysis, were mostly used as helpful tools in the early design stage. Their real-world effect was not as significant when compared to the primary UML artifacts.

For upcoming projects, it's a good idea to keep hold of functional and non-functional requirements, user stories, activity diagrams, and class diagrams. These practices help to clearly define what needs to be done and what the goals are, keep the focus on the user, allow for an early look at how the system will work, and create a solid base for the implementation.



% --------------------------------
\section{Refactoring}

How did refactoring help your project's code quality, skills balance in
the group, etc.? (Delete this section if refactoring was not used.)
% --------------------------------
\section{Code Review}

How did code review help with code quality and knowledge sharing?
What were the challenges and how did you overcome them?

% --------------------------------
\section{Testing}

To what extent did testing help the quality of your product? What were
the disadvantages — e.g.\ slow productivity, intellectually unstimulating
— and how do you plan to adopt testing in future projects?

% --------------------------------
\section{Overall Group Reflection}

Include an overall reflection on the project outcome:

\begin{itemize}
    \item Whether all requirements were met or some were changed
          (e.g.\ a comparison table).
    \item Whether the product solves the stated issue.
    \item Whether the product leads to unforeseen problems that would
          be addressed with future development.
    \item Possible future development in terms of features, scalability,
          or other important factors.
\end{itemize}